\section{Computing an active odd-cut constraint}
\label{d:sec-odd-cuts}

The optimization in Section~\ref{d:sec-optimization} needs both the value of
\(w(x)\) and an affine function attaining its defining maximum.
The degree violations are immediate to compute.  We supply the remaining
ingredient: an exact minimum odd cut, together with its vertex set.
The use of minimum odd cuts for matching-constraint separation goes back
to Padberg and Rao~\citep{PadbergRao1982}. We prove the recursion and its
arithmetic bounds here.

Throughout this section, an auxiliary graph has vertex set \(W\) and
nonnegative integer capacities \(c_{ij}=c_{ji}\) for distinct \(i,j\in W\).
Absent edges have capacity zero.  Parallel capacities are summed and loops
are discarded.  Write
\[
 c(S)=\sum_{\substack{i\in S\\ j\in W\setminus S}}c_{ij},
 \qquad
 C=\sum_{\{i,j\}\subseteq W}c_{ij},
 \qquad v=|W|.
\]

We use the classical augmenting-path method of Ford and Fulkerson
\citep[Section~2]{FordFulkerson1957}, specialized here to unit
augmentations. The capacity-dependent bound needed below is proved
explicitly.

\begin{lemma}[An integral two-terminal minimum cut]
\label{d:lem-unit-mincut}
For distinct \(a,b\in W\), a minimum cut separating \(a\) from \(b\)
can be found deterministically using at most \(C\) unit augmentations
and \(O(v^2(C+1))\) arithmetic operations on integers of
\(O(\log_2(v+C+2))\) bits.
\end{lemma}

\begin{proof}
Maintain integers \(b_{ij}=-b_{ji}\) such that
\(b_{ij}\le c_{ij}\) in both directions and
\(\sum_j b_{ij}=0\) for \(i\notin\{a,b\}\).
The two directional inequalities give
\(-c_{ij}\le b_{ij}\le c_{ij}\).
Start with zero flow.  If a simple directed path from \(a\) to \(b\)
uses only steps with residual capacity \(c_{ij}-b_{ij}>0\), increase
the flow on each forward step by one and decrease its reverse by one.
Integrality makes this feasible, conservation is preserved, and the
net source outflow \(q=\sum_j b_{aj}\) increases by one.
Since
\[
 q\le \sum_j c_{aj}\le C,
\]
there are at most \(C\) augmentations.

When no such path exists, let \(S\) be the vertices reachable from \(a\)
by positive-residual steps.  Then \(b\notin S\), and every step from
\(S\) to its complement is saturated.  Antisymmetry cancels internal
flows, so conservation gives
\[
 q=\sum_{\substack{i\in S\\j\notin S}}b_{ij}
   =\sum_{\substack{i\in S\\j\notin S}}c_{ij}=c(S).
\]
For every other \(a\)--\(b\) cut, the same conservation identity and
\(b_{ij}\le c_{ij}\) give capacity at least \(q\).  Thus \(S\) is minimum.
A graph search with fixed vertex order finds a simple residual path,
or the final reachable set, in \(O(v^2)\) operations.  There are at most
\(C+1\) searches, and all flows and residual capacities have magnitude
at most \(2C\).  The argument includes \(C=0\).
\end{proof}

\begin{lemma}[Minimum odd-terminal cut]
\label{d:lem-odd-cut-oracle}
Let \(U\subseteq W\) have positive even cardinality.  One can
deterministically find a set \(S\subseteq W\) minimizing \(c(S)\)
subject to \(|S\cap U|\) being odd.  The algorithm returns the set itself,
uses at most \(|U|-1\) recursive instances, and requires
\(O(|U|^3v^2(C+1))\) arithmetic operations on integers of
\(O(\log_2(v+C+2))\) bits.
\end{lemma}

\begin{proof}
Call a cut \emph{terminal-separating} if both sides meet \(U\).
Compute a minimum cut for every pair of distinct terminals using
Lemma~\ref{d:lem-unit-mincut}, and take the least-capacity result,
with sides \(A\) and \(\overline A=W\setminus A\).
This is a minimum terminal-separating cut: every such cut separates
some terminal pair, and every recorded pair cut is terminal-separating.

If \(|A\cap U|\) is odd, return \(A\).  Every admissible odd-terminal
cut is terminal-separating, so this answer is optimal.
Otherwise both \(A\cap U\) and \(\overline A\cap U\) have positive even
cardinality.  Form two children:
\begin{enumerate}[label=(\roman*),nosep]
\item keep \(A\), contract \(\overline A\) to a new nonterminal, and
      use terminal set \(A\cap U\);
\item keep \(\overline A\), contract \(A\) to a new nonterminal, and
      use terminal set \(\overline A\cap U\).
\end{enumerate}
For example, in the first child the capacity from \(i\in A\) to the
new vertex is \(\sum_{j\in\overline A}c_{ij}\).
Recurse on both children and return the better lifted cut.
Every call strictly decreases the positive even terminal count;
when \(|U|=2\), the first cut is already odd.

\paragraph{Lifting cuts.}
A child cut lifts by replacing its contracted vertex, if included,
with the entire block it represents.  The crossing capacities are
unchanged, and including that block adds an even number of parent
terminals.  Thus the lift has the same capacity and is terminal-odd.

This remains true through arbitrarily nested contractions.
Relative to the initial instance, current vertices form a partition
into blocks.  Every current terminal is an original terminal singleton,
and every nonterminal block contains an even number of original terminals.
Initially the latter count is zero; each contraction combines an even
number of current terminals with blocks of even parity, preserving
the assertion.  Capacities between current vertices are sums of
original capacities between their blocks.
Consequently every returned subset lifts to an original subset with
the same capacity and terminal parity.
Figure~\ref{d:fig-odd-cut-lift} illustrates the parity rule.

\paragraph{Preserving an optimum.}
Let \(S\) be an optimal odd-terminal side in the current instance.
Exactly one of \(S\cap A\) and \(S\cap\overline A\) contains an odd number
of terminals; exchange the names of the two sides so that it is \(S\cap A\).
Because \(|A\cap U|\) is even, replacing \(S\) by its complement
preserves this odd intersection.  Make this replacement if \(S\)
contains all terminals of \(\overline A\).
Now \(S\cup A\) contains terminals of \(A\) and omits a terminal of
\(\overline A\), so
\[
 c(S\cup A)\ge c(A).
\]
For any two vertex sets, nonnegative undirected capacities satisfy
\begin{equation}
\label{d:eq-cut-submodularity}
 c(S\cap A)+c(S\cup A)\le c(S)+c(A).
\end{equation}
Indeed, the right side minus the left side is twice the capacity
between \(S\setminus A\) and \(A\setminus S\).
It follows that \(c(S\cap A)\le c(S)\).
The set \(S\cap A\) is a feasible odd-terminal side in the first child,
avoiding its contracted vertex.  Hence at least one child has optimum
at most the parent's.  Conversely, lifting shows that every child
optimum is at least the parent's.  Induction on \(|U|\) proves that
the better recursive answer attains the parent optimum.

\paragraph{Size and running time.}
Each nonleaf splits its terminal set into two disjoint nonempty even
sets.  There are at most \(|U|/2\) leaves and hence at most \(|U|-1\)
instances in this binary recursion tree.
Each contracted side contains at least two terminals, so a child
also has fewer vertices than its parent.
Within any one instance, every surviving capacity sums a disjoint
set of original edge capacities; original edges internal to blocks
are discarded.  Its total capacity is therefore at most \(C\).
Capacities in different children may overlap, and we account for them
by multiplying the cost of an instance by the number of instances.

There are at most \(O(|U|^2)\) pair computations per instance,
each costing \(O(v^2(C+1))\) by Lemma~\ref{d:lem-unit-mincut}.
Contractions and lifted-set bookkeeping fit within the same bound.
This proves the stated running time.
All comparisons and tie choices may use fixed orderings; zero
capacities require no perturbation.
\end{proof}

\begin{figure}[tb]
\centering
\begin{tikzpicture}[
  x=1cm,y=1cm,
  terminal/.style={circle,fill=black,inner sep=2.3pt},
  contracted/.style={circle,draw=black,fill=white,line width=.7pt,inner sep=3pt},
  block/.style={draw=black!55,rounded corners=3pt},
  cutshade/.style={fill=blue!12,draw=none},
  edge/.style={draw=black!50,line width=.55pt},
  every node/.style={font=\small}
]
\path[cutshade] (0.25,1.3)--(2.55,1.3)--(2.55,.2)--(4.55,.2)
               --(4.55,2.25)--(.25,2.25)--cycle;
\draw[block] (.15,.1) rectangle (2.25,2.35);
\draw[block] (2.45,.1) rectangle (4.65,2.35);
\node at (1.2,2.68) {$A$};
\node at (3.55,2.68) {$\overline A$};
\node[terminal] (a1) at (1.2,1.8) {};
\node[terminal] (a2) at (1.2,.65) {};
\node[terminal] (b1) at (3.55,1.8) {};
\node[terminal] (b2) at (3.55,.65) {};
\draw[edge] (a1)--(a2);
\draw[edge] (a1)--(b1);
\draw[edge] (a1)--(b2);
\draw[edge] (a2)--(b1);
\draw[edge] (a2)--(b2);
\draw[edge] (b1)--(b2);
\node at (2.4,-.3) {lifted side: three terminals};

\draw[-{Stealth[length=2mm]},line width=.7pt]
  (4.95,1.2)--(6.35,1.2);
\node[align=center,font=\footnotesize] at (5.65,1.8)
  {contract\\$\overline A$};

\path[cutshade] (6.7,1.3) rectangle (10.8,2.25);
\draw[block] (6.6,.1) rectangle (8.7,2.35);
\node at (7.65,2.68) {$A$};
\node[terminal] (c1) at (7.65,1.8) {};
\node[terminal] (c2) at (7.65,.65) {};
\node[contracted] (z) at (10.2,1.8) {};
\draw[edge] (c1)--(c2);
\draw[edge] (c1)--(z);
\draw[edge] (c2)--(z);
\node[align=center,font=\footnotesize] at (10.2,.7)
  {nonterminal\\even number of\\original terminals};
\node at (8.7,-.3) {child side: one terminal};
\end{tikzpicture}
\caption{Lifting through an even-terminal contraction.
Filled dots are terminals and the open dot represents the contracted
block.  In this schematic instance each side initially has two terminals.
The shaded child side lifts to the shaded parent side by taking the
entire contracted block, preserving odd parity.  Parallel crossing
capacities are added, so the cut capacity is preserved as well.}
\label{d:fig-odd-cut-lift}
\end{figure}


\paragraph{Recovering an active row in the original coordinates.}
At an optimization trial point, \(D_0x_e\) is a nonnegative integer
for every original complete-graph pair \(e\).
Indeed, the initial point is zero, and an integral supergradient step
of size \(1/D_0\), followed by clipping to \([0,1]\), preserves this grid.
Apply Lemma~\ref{d:lem-odd-cut-oracle} with \(U=V\) and capacities
\(c_e=D_0x_e\).  Store each contracted vertex's block of original
vertices, so the returned cut yields an explicit original set
\(S\subseteq V\) of odd cardinality.  It minimizes \(x(\delta(S))\).

The corresponding affine function in the definition of \(w\) is
\[
 \ell_S(y)=1-\sum_{e\in\delta(S)}y_e.
\]
Its coefficient is \(-1\) on \emph{every} original pair crossing \(S\),
including pairs with \(x_e=0\), and zero elsewhere.
Thus the row is recovered by scanning the original pairs, even if
zero-capacity edges were omitted during graph searches.
Edges discarded as contraction loops have both endpoints in one
block and do not cross a lifted cut.
Comparing \(\ell_S(x)\) with zero and the two signed degree violations
for every vertex gives \(w(x)\) and an affine row attaining it, as
required in Section~\ref{d:sec-optimization}.

Here the initial total capacity is at most \(pD_0\).
A merged edge can exceed \(D_0\), but the bound on each instance's
total capacity remains \(pD_0\); likewise a large reverse residual
does not change the source-outflow bound in Lemma~\ref{d:lem-unit-mincut}.
One complete violation computation therefore takes
\[
 O\!\left(n^5(pD_0+1)\right)
\]
arithmetic operations.
The numerical bounds on \(D_0\) and the resulting bit complexity
are recorded in Proposition~\ref{d:prop-bit-complexity}.
